% ============================================================
%  Introduzione alla Parte I: Preliminari
% ============================================================

\chapter*{Introduzione Parte I}
\addcontentsline{toc}{chapter}{Introduzione alla Parte~I}

I quattro capitoli di questa parte raccolgono gli strumenti 
matematici e concettuali su cui poggia l'intera trattazione, 
insieme alle convenzioni notazionali adottate sistematicamente 
nel testo. Non sono un corso autonomo di algebra lineare o di 
teoria vettoriale: sono una selezione mirata degli strumenti 
che il testo utilizza sistematicamente, presentati con il 
livello di dettaglio e con l'enfasi che il loro impiego nella 
meccanica delle strutture richiede. Lo studente che ha già acquisito 
questi strumenti in corsi precedenti può consultarli come 
riferimento; chi li incontra per la prima volta troverà qui 
una trattazione autosufficiente, orientata fin dall'inizio 
verso le applicazioni strutturali.

\section*{Il programma della parte}

Il \textbf{Capitolo~\ref{app:notazioni}} fissa le 
\emph{notazioni e le convenzioni} di uso ricorrente nel testo: 
la distinzione fra punti dello spazio euclideo, vettori 
geometrici e vettori numerici, applicazioni lineari e matrici 
rappresentative, con le rispettive convenzioni tipografiche; 
e il significato di simboli matematici di impiego sistematico, 
quali il segno di eguaglianza nei suoi quattro significati, le 
distinzioni fra costanti, parametri e variabili, fra salto e 
incremento, fra costante (nel tempo) e uniforme (nello spazio). 
Le convenzioni specialistiche dei modelli strutturali sono 
invece introdotte nei capitoli in cui se ne fa primo uso.

Il \textbf{Capitolo~\ref{app:algebra_eqlineari}} richiama 
l'\emph{algebra delle matrici e i sistemi di equazioni lineari}. 
La formulazione e la soluzione dei problemi strutturali si 
appoggiano sistematicamente sul linguaggio dell'algebra lineare: 
vettori numerici, matrici, operazioni fondamentali, 
classificazione e risoluzione dei sistemi lineari. Particolare 
attenzione è dedicata alla distinzione fra i quattro casi --- 
isodeterminato, sottodeterminato, sovradeterminato, degenere --- 
che si ritrovano puntualmente nella classificazione cinematica e 
statica dei sistemi strutturali.

Il \textbf{Capitolo~\ref{app:applicazionilin_dualita}} sviluppa 
la geometria delle \emph{applicazioni lineari e la dualità}. 
I quattro sottospazi fondamentali --- nucleo, immagine, nucleo 
aggiunto, immagine aggiunta --- forniscono la chiave geometrica 
per interpretare la classificazione dei sistemi lineari del 
capitolo precedente e, soprattutto, per comprendere la struttura 
profonda del problema meccanico: l'operatore cinematico e il 
suo aggiunto, l'operatore di equilibrio, la dualità fra 
spostamenti e forze. L'identità bilineare che emerge da questa 
analisi è la forma astratta del principio dei lavori virtuali, 
che costituisce il filo conduttore dell'intera opera.

Il \textbf{Capitolo~\ref{app:vettori_geometrici}} presenta i 
\emph{vettori geometrici}: liberi, applicati, polari e assiali. 
Le forze, gli spostamenti, i momenti e le rotazioni sono vettori 
geometrici di tipo diverso, e la distinzione fra vettori polari 
e vettori assiali --- fondata sul diverso comportamento sotto 
riflessioni --- ha conseguenze fisiche precise che si riflettono 
nella rappresentazione grafica e nel calcolo. Il momento di una 
forza rispetto a un polo, la sua legge di variazione al variare 
del polo e l'invariante scalare che ne emerge sono strumenti 
che si ritrovano identici nella statica del corpo rigido e 
nella cinematica degli spostamenti rigidi --- prima 
manifestazione della dualità cinematica-statica.

\section*{Uno strumento, non un fine}

È opportuno dichiarare esplicitamente il ruolo di questa parte. 
I preliminari non sono un fine in sé: sono uno strumento. Il 
loro scopo è fornire al lettore il linguaggio e i concetti 
necessari per seguire la trattazione meccanica senza 
interruzioni. Per questa ragione la presentazione è orientata 
verso le applicazioni: ogni concetto è introdotto con un occhio 
al suo impiego nella meccanica delle strutture, e i risultati più 
astratti sono accompagnati da anticipazioni del contesto in 
cui appariranno.

In particolare, il Capitolo~\ref{app:applicazionilin_dualita} 
sulla dualità non è uno strumento ausiliario ma una componente 
essenziale dell'architettura concettuale del testo. Chi lo 
attraversa con attenzione troverà, nei capitoli successivi, 
una teoria strutturale in cui ogni elemento ha il suo posto 
preciso e le connessioni fra le parti sono trasparenti.